\section{AION Phase 20Y --- Canonical Payload Hashing}
\label{sec:aion-phase20y-canonical-payload-hashing}

\subsection{Status}

Phase 20Y locks canonical checkpoint payload hashing.

\begin{quote}
\textbf{Status: LOCKED --- checkpoint payload identity is canonical, deterministic, and recalculated immediately before execution.}
\end{quote}

\subsection{Purpose}

Checkpoint approval is only safe if the human approves the exact payload that will later be executed.

Phase 20Y ensures that checkpoint payloads are represented by deterministic canonical hashes and that any post-approval edit invalidates the approval.

\begin{quote}
\textbf{An approval belongs to one exact payload hash only.}
\end{quote}

\subsection{Canonical JSON Rule}

Checkpoint payloads MUST be canonicalised using:

\begin{verbatim}
json.dumps(payload, sort_keys=True, separators=(",", ":"), ensure_ascii=False)
\end{verbatim}

This removes non-semantic key-order and whitespace differences from the hash calculation.

\subsection{Payload Hash}

The payload hash MUST be:

\begin{verbatim}
sha256:<digest>
\end{verbatim}

computed over canonical JSON.

\subsection{Payload Hash Record}

A payload hash record MUST include:

\begin{itemize}
  \item mission ID;
  \item mission run ID;
  \item checkpoint ID;
  \item step ID;
  \item action type;
  \item payload hash;
  \item canonical payload;
  \item hash algorithm;
  \item schema version.
\end{itemize}

\subsection{Execution-Time Recalculation}

Immediately before execution, the runtime MUST recalculate the payload hash from the current payload.

If:

\begin{equation}
H_{\mathrm{approved\_payload}}
\neq
H_{\mathrm{current\_payload}}
\end{equation}

then execution MUST be blocked.

\subsection{Edited Payload Invalidation}

If a payload changes after approval, the runtime MUST return:

\begin{verbatim}
approval_valid: false
requires_fresh_approval: true
reason: payload_changed_after_approval
\end{verbatim}

\subsection{Validation Hash}

Every execution-time payload validation MUST emit a deterministic validation hash binding:

\begin{itemize}
  \item approved payload hash;
  \item current payload hash;
  \item payload unchanged decision;
  \item execution allowed decision;
  \item reason.
\end{itemize}

\subsection{Validation}

\begin{verbatim}
python -m pytest -q \
  backend/tests/workflow_capsules/test_aion_phase20a_mission_contract_lock.py \
  backend/tests/workflow_capsules/test_aion_phase20b_mission_template_lock.py \
  backend/tests/workflow_capsules/test_aion_phase20c_mission_planner_lock.py \
  backend/tests/workflow_capsules/test_aion_phase20d_autonomy_lane_classifier_lock.py \
  backend/tests/workflow_capsules/test_aion_phase20v_kernel_guard_lock.py \
  backend/tests/workflow_capsules/test_aion_phase20e_mission_runtime_lock.py \
  backend/tests/workflow_capsules/test_aion_phase20aa_business_container_artifacts_lock.py \
  backend/tests/workflow_capsules/test_aion_phase20z_pilot_native_runtime_lock.py \
  backend/tests/workflow_capsules/test_aion_phase20w_self_repair_runtime_lock.py \
  backend/tests/workflow_capsules/test_aion_phase20x_session_vfs_lock.py \
  backend/tests/workflow_capsules/test_aion_phase20f_mission_approval_lock.py \
  backend/tests/workflow_capsules/test_aion_phase20y_canonical_payload_hashing_lock.py
\end{verbatim}

\subsection{Lock Footer}

\begin{verbatim}
Lock ID: AION-PHASE20Y-CANONICAL-PAYLOAD-HASHING
Status: LOCKED
Maintainer: Tessaris AI
Author: Kevin Robinson
\end{verbatim}

\subsection{Phase 20Y.1 --- Numeric Primitive Normalisation and Validation Hash Closure}
\label{subsec:aion-phase20y1-numeric-normalisation-validation-hash}

\subsubsection{Status}

Phase 20Y.1 hardens canonical payload hashing against primitive representation drift.

\begin{quote}
\textbf{Status: LOCKED --- payload primitives MUST be normalised before canonical hashing, and validation hashes MUST bind both approved and current payload states.}
\end{quote}

\subsubsection{Numeric Typing and Encoding Normalisation Invariant}

To prevent non-semantic floating-point representation variance, integer-to-float promotion, or Unicode escape manipulation from causing false approval invalidation or semantic bypassing, the payload engine MUST run a data normaliser before serialization.

The normaliser MUST apply recursively across the payload object graph:

\begin{itemize}
  \item dictionaries are traversed deterministically;
  \item lists are traversed recursively;
  \item strings are normalised to canonical Unicode form;
  \item booleans and null values are preserved;
  \item integers are preserved;
  \item integral floats such as \texttt{100.0} are normalised to \texttt{100};
  \item non-integral floats are normalised to a fixed precision bound.
\end{itemize}

Formally:

\begin{equation}
\mathrm{NormalizePrimitives}(Payload)
\longrightarrow
Payload_{\mathrm{normalised}}
\end{equation}

The compact JSON serializer MUST consume only the normalised structural layout.

\subsubsection{Primitive Representation Drift}

The following payloads MUST hash identically after normalisation:

\begin{verbatim}
{"max_limit": 100}
{"max_limit": 100.0}
\end{verbatim}

The following Unicode-equivalent strings MUST hash identically after normalisation:

\begin{verbatim}
"Café"
"Cafe\u0301"
\end{verbatim}

\subsubsection{Validation Hash Formula Closure}

The validation hash acts as a proof record that execution-time validation occurred before an operational block was allowed or blocked.

The validation hash is defined as:

\begin{equation}
\mathcal{V}_{\mathrm{hash}}
=
\mathrm{SHA256}
(
H_{\mathrm{approved\_payload}}
\parallel
H_{\mathrm{current\_payload}}
\parallel
\mathrm{DecisionFlags}
\parallel
\texttt{StepID}
)
\end{equation}

The decision flags include:

\begin{itemize}
  \item \texttt{payload\_unchanged}
  \item \texttt{execution\_allowed}
  \item \texttt{reason}
\end{itemize}

\subsubsection{Execution Boundary Rule}

Immediately before a checkpointed action executes, the runtime MUST:

\begin{enumerate}
  \item normalise the current payload primitives;
  \item canonicalise the normalised payload;
  \item calculate the current payload hash;
  \item compare it to the approved payload hash;
  \item emit the validation hash closure;
  \item block execution if the hashes differ.
\end{enumerate}

\subsubsection{Validation Coverage}

Phase 20Y.1 tests prove:

\begin{itemize}
  \item integral floats and integers hash identically;
  \item floating precision is deterministic;
  \item Unicode-equivalent strings hash identically;
  \item normalised execution validation allows equivalent integer/float payloads;
  \item validation hash closure changes when decision state changes.
\end{itemize}

\subsubsection{Security Invariant}

\begin{quote}
\textbf{AION Pilot checkpoint payload hashing is based on canonical structure, not unstable runtime formatting. Meaningful edits invalidate approval; non-semantic primitive representation drift does not.}
\end{quote}

\subsubsection{Lock Footer}

\begin{verbatim}
Lock ID: AION-PHASE20Y1-NUMERIC-PRIMITIVE-NORMALISATION-VALIDATION-HASH-CLOSURE
Status: LOCKED
Maintainer: Tessaris AI
Author: Kevin Robinson
\end{verbatim}
